\chapter{Introduction}

\section{Overview}
Aviation safety relies heavily on passengers understanding emergency protocols. However, traditional methods of information delivery are passive. Passengers often ignore safety cards and videos, leading to a lack of preparedness in real emergencies. \textbf{Virtual Reality (VR)} offers a solution by creating a "presence" that traditional media cannot match. This project aims to transform the safety briefing from a monologue into an interactive experience. By simulating the cabin environment and requiring user input to progress, the system ensures that the passenger is mentally present and actively processing the safety information.

\section{Motivation}
The primary motivation for this project is to address the critical gap in passenger safety education. Studies have consistently shown that passengers pay little attention to pre-flight safety briefings. In an emergency, this lack of knowledge can be fatal. VR technology provides a unique opportunity to simulate emergency scenarios in a safe, controlled environment, allowing passengers to practice procedures such as donning oxygen masks or locating emergency exits.

\section{Problem Definition}
The core problem is the ineffectiveness of current safety briefing methods. Passengers are passive recipients of information, which leads to low retention rates. The challenge is to create a system that engages the user actively and verifies their understanding of safety protocols.

\section{Organization of the Report}
The remainder of this report is organized as follows:
\begin{itemize}
    \item Chapter 2 presents a survey of existing literature and related work.
    \item Chapter 3 defines the problem statement and project objectives.
    \item Chapter 4 details the software requirements and specifications.
    \item Chapter 5 describes the system design and architecture.
    \item Chapter 6 outlines the proposed methodology.
    \item Chapter 7 discusses the implementation details.
    \item Finally, the report concludes with references.
\end{itemize}
